% =====================================================================
%  CHƯƠNG 1 — LỊCH SỬ PHÁT TRIỂN TRÍ TUỆ NHÂN TẠO
% =====================================================================
\chapter{Lịch Sử Phát Triển Trí Tuệ Nhân Tạo}
\label{ch:lich-su}

\section{Tổng quan bức tranh hơn tám thập kỷ của AI}

Lịch sử trí tuệ nhân tạo không phải một đường đi lên đều đặn. Đó là chuỗi các chu kỳ kỳ vọng: mỗi đột phá kéo theo
lời hứa vượt quá năng lực thật của công nghệ, rồi một giai đoạn thoái trào khi tiền tài trợ cạn, gọi là
\emph{mùa đông AI}. Qua các chu kỳ đó, cách máy tính “suy nghĩ” đổi ba lần: từ suy diễn trên ký hiệu do con người
viết, sang mô hình xác suất học từ dữ liệu, rồi tới học biểu diễn nhiều tầng của học sâu. Hình~\ref{fig:timeline}
đặt mười sáu mốc then chốt lên cùng một trục thời gian, kèm các kỷ nguyên tương ứng.

\hinhtikz{f1_timeline}{Dòng thời gian tiến hóa của trí tuệ nhân tạo (1943 -- 2026): các mốc then chốt và các kỷ
nguyên công nghệ. Hai dải xám là hai mùa đông AI (1974 -- 1980 và 1987 -- 1993), dải cam giữa chúng là thời kỳ hệ chuyên
gia (1980 -- 1987).}{timeline}

Hiểu lịch sử không chỉ để biết ai làm gì, năm nào. Mỗi mô hình được hiện thực ở các chương sau đều là lời đáp cho
một thất bại cụ thể trong quá khứ: MLP và lan truyền ngược trả lời giới hạn XOR của Perceptron, SVM và Random Forest
trả lời sự giòn của hệ chuyên gia, LSTM trả lời hiện tượng triệt tiêu đạo hàm của RNN, ResNet trả lời sự thoái hóa của
mạng rất sâu.

\section{Khởi nguồn và kỷ nguyên logic biểu tượng (1940s -- 1960s)}

\subsection{Ý niệm về máy tính có tư duy và phép thử Turing}

Năm 1950, Alan Turing công bố bài báo \emph{“Computing Machinery and Intelligence”} trên tạp chí \emph{Mind}
\cite{turing1950}, mở đầu bằng câu hỏi “Máy móc có thể suy nghĩ không?”. Để tránh tranh cãi triết học về khái niệm
“ý thức”, Turing thay câu hỏi bằng một thí nghiệm hành vi, \emph{trò chơi bắt chước}: một người thẩm vấn trò chuyện
bằng văn bản với hai đối tượng giấu mặt, một người và một máy. Nếu người thẩm vấn không phân biệt được đâu là máy,
ta nói máy đã thể hiện hành vi thông minh. Phép thử Turing là phương pháp luận đầu tiên định lượng hóa trí tuệ nhân
tạo bằng hành vi quan sát được, và đến nay vẫn là điểm quy chiếu khi bàn về các mô hình ngôn ngữ lớn.

\subsection{Mô hình nơ-ron hình thức và Hội thảo Dartmouth (1956)}

Trước cả Turing, năm 1943, nhà thần kinh học Warren McCulloch và nhà logic học Walter Pitts đề xuất mô hình
nơ-ron hình thức \cite{mcculloch1943}: một nơ-ron tính tổng có trọng số các tín hiệu nhị phân rồi so với một ngưỡng,
\begin{equation}
  y = f\Big(\sum_{i=1}^{n} w_i x_i - \theta\Big), \qquad x_i \in \{0, 1\},
  \label{eq:mcp}
\end{equation}
với $f$ là hàm bước Heaviside. Hai tác giả chứng minh mạng các nơ-ron như vậy thực hiện được mọi hàm logic AND, OR,
NOT. Năm 1949, Donald Hebb bổ sung quy tắc học “các nơ-ron cùng phát xung thì liên kết mạnh lên”, cơ sở sinh học đầu
tiên cho việc \emph{thay đổi} trọng số.

Ngành AI chính thức ra đời mùa hè năm 1956 tại Hội thảo Dartmouth do John McCarthy, Marvin Minsky, Claude Shannon và
Nathaniel Rochester tổ chức. McCarthy đặt tên “Artificial Intelligence” với giả thuyết: mọi khía cạnh của việc học hay
của trí tuệ, về nguyên tắc, đều mô tả được chính xác tới mức một cỗ máy có thể mô phỏng nó.

\subsection{Perceptron của Rosenblatt và các hệ suy diễn ký hiệu đầu tiên}

Năm 1958, Frank Rosenblatt công bố Perceptron \cite{rosenblatt1958}, mô hình nơ-ron đầu tiên \emph{tự học} trọng số
từ dữ liệu có nhãn theo quy tắc
\begin{equation}
  \mathbf{w}^{(t+1)} = \mathbf{w}^{(t)} + \eta\,(y - \hat y)\,\mathbf{x},
  \label{eq:perceptron}
\end{equation}
với $\eta > 0$ là tốc độ học. Định lý hội tụ Perceptron bảo đảm thuật toán dừng sau hữu hạn bước nếu dữ liệu khả
tách tuyến tính. Song song với hướng kết nối (connectionism) này, hướng biểu tượng (symbolic AI) cũng gặt hái thành
công: chương trình Logic Theorist (1956) của Newell, Simon và Shaw chứng minh được 38 trên 52 định lý ở chương 2 của
\emph{Principia Mathematica}; General Problem Solver (1959) giải các bài toán đố bằng phân tích phương tiện -- mục đích.

\section{“Mùa đông AI” lần thứ nhất (1970s)}

\subsection{Giới hạn toán học của Perceptron và cuốn sách của Minsky \& Papert (1969)}

Năm 1969, Minsky và Papert xuất bản \emph{Perceptrons} \cite{minsky1969}, chứng minh chặt chẽ rằng Perceptron một tầng
chỉ học được các hàm khả tách tuyến tính, và vì vậy thất bại với cả hàm XOR đơn giản (Bảng~\ref{tab:xor}).

\begin{table}[H]\centering
\caption{Bảng chân lý của hàm XOR: hai lớp nằm trên hai đường chéo nên không đường thẳng nào tách được.}\label{tab:xor}
\begin{tabular}{cccl}
\toprule
$x_1$ & $x_2$ & $y = x_1 \oplus x_2$ & Vị trí trên mặt phẳng $\mathbb{R}^2$ \\
\midrule
0 & 0 & 0 & gốc tọa độ $(0, 0)$ -- nhãn âm \\
0 & 1 & 1 & điểm $(0, 1)$ -- nhãn dương \\
1 & 0 & 1 & điểm $(1, 0)$ -- nhãn dương \\
1 & 1 & 0 & điểm $(1, 1)$ -- nhãn âm \\
\bottomrule
\end{tabular}
\end{table}

Muốn tách $\{(0,1), (1,0)\}$ khỏi $\{(0,0), (1,1)\}$ cần một đường thẳng $w_1x_1 + w_2x_2 + b = 0$ để hai nhóm nằm hai
phía, nhưng từ $b < 0$, $w_1 + b > 0$, $w_2 + b > 0$ suy ra $w_1 + w_2 + b > -b > 0$, mâu thuẫn với yêu cầu điểm
$(1,1)$ phải âm. Minsky và Papert biết mạng nhiều tầng giải được XOR, nhưng thời đó chưa ai có cách huấn luyện các
tầng ẩn. Cuốn sách khiến nguồn tài trợ cho hướng nơ-ron gần như đóng băng suốt hơn một thập kỷ.

\subsection{Báo cáo Lighthill (1973) và sự cắt giảm tài trợ của DARPA}

Năm 1973, Báo cáo Lighthill \cite{lighthill1973} gửi Hội đồng Nghiên cứu Khoa học Anh kết luận rằng các hệ AI ký hiệu
chỉ chạy được trong “thế giới đồ chơi” và bất lực trước bài toán thật vì \emph{bùng nổ tổ hợp}: số trạng thái cần tìm
kiếm tăng theo hàm mũ của kích thước bài toán. Chính phủ Anh cắt tài trợ AI ở hầu hết các trường đại học; cùng lúc, tu
chính án Mansfield khiến DARPA của Hoa Kỳ ngừng chi cho nghiên cứu không có ứng dụng quân sự trực tiếp. Giai đoạn
1974 -- 1980 được gọi là mùa đông AI thứ nhất.

\section{Sự trỗi dậy của hệ chuyên gia và “mùa đông AI” lần thứ hai (1980s -- 1990s)}

\subsection{Kỷ nguyên vàng của hệ chuyên gia}

Đầu thập niên 1980, AI hồi sinh khi chuyển từ tìm kiếm tổng quát sang \emph{tri thức chuyên ngành hẹp}. Một hệ chuyên
gia tách hai thành phần: \textbf{cơ sở tri thức} gồm sự kiện và luật dạng \texttt{IF} điều kiện \texttt{THEN} kết luận,
và \textbf{máy suy diễn} chạy luật theo chiều tiến (từ dữ kiện ra kết luận) hoặc chiều lùi (từ giả thuyết truy về dữ
kiện). Các hệ tiêu biểu gồm DENDRAL (xác định cấu trúc phân tử từ phổ khối), MYCIN (chẩn đoán nhiễm trùng máu, độ chính
xác ngang bác sĩ chuyên khoa) và XCON/R1 của Digital Equipment Corporation, tự cấu hình máy VAX và được ước tính tiết
kiệm khoảng 40 triệu USD mỗi năm. Thành công của XCON kéo theo làn sóng đầu tư vào máy Lisp chuyên dụng và dự án Máy
tính Thế hệ Thứ năm của Nhật Bản.

\subsection{Sự tái sinh của mạng nơ-ron: thuật toán lan truyền ngược (1986)}

Năm 1986, Rumelhart, Hinton và Williams công bố trên \emph{Nature} thuật toán lan truyền ngược cho mạng nhiều tầng
\cite{rumelhart1986}. Dùng quy tắc dây chuyền với hàm kích hoạt khả vi như sigmoid $\sigma(z) = 1/(1+e^{-z})$, đạo hàm
của mất mát $L$ theo một trọng số bất kỳ ở tầng ẩn $l$ tách thành tích của sai số truyền ngược và kích hoạt tầng trước:
\begin{equation}
  \frac{\partial L}{\partial w^{(l)}_{ij}} = \delta^{(l)}_i \, a^{(l-1)}_j .
  \label{eq:bp}
\end{equation}
Kết quả này phá “lời nguyền XOR”: mạng hai tầng học được biểu diễn ẩn phi tuyến. Cùng giai đoạn, Fukushima
(1980) đề xuất Neocognitron \cite{fukushima1980} với các tầng trích đặc trưng cục bộ (S-cells) xen kẽ các tầng bất biến
vị trí (C-cells), tiền thân trực tiếp của CNN. Chương~\ref{ch:hoc-may} hiện thực lại công thức~\eqref{eq:bp} bằng NumPy.

\subsection{“Mùa đông AI” lần thứ hai (1987 -- 1993)}

Hệ chuyên gia sớm lộ ba điểm yếu cấu trúc. \emph{Nút thắt thu nhận tri thức}: phỏng vấn chuyên gia để mã hóa hàng vạn
luật vừa đắt vừa chậm, lại sinh mâu thuẫn. \emph{Tính giòn}: hệ thống trả lời sai hoặc lặp vô tận khi gặp tình huống
nằm ngoài các luật đã viết, chi phí bảo trì vượt lợi ích. \emph{Sụp đổ phần cứng}: máy trạm Unix và vi xử lý x86 rẻ
hơn khiến máy Lisp lỗi thời, kéo theo phá sản dây chuyền. Hình~\ref{fig:hype} minh họa định tính các chu kỳ kỳ vọng và
hai mùa đông AI.

\hinhtikz{f1_hype}{Chu kỳ kỳ vọng và hai thời kỳ “mùa đông AI” (minh họa định tính, không phải số đo).}{hype}

Sau mỗi lần thoái trào, cộng đồng tái cấu trúc phương pháp luận theo hướng thực dụng hơn và có cơ sở toán học chắc hơn.

\section{Kỷ nguyên học máy thống kê (1990s -- 2010)}

Bài học rút ra từ thất bại của hệ chuyên gia: thay vì cố viết tay mọi quy tắc của thế giới, hãy để máy học quy luật
xác suất trực tiếp từ dữ liệu.

\subsection{Lý thuyết học thống kê và máy vector hỗ trợ (SVM)}

Vapnik và Chervonenkis xây dựng lý thuyết học thống kê với nguyên lý cực tiểu hóa rủi ro cấu trúc: kiểm soát cùng lúc
sai số huấn luyện và độ phức tạp của lớp hàm (đo bằng chiều VC). Đỉnh cao là máy vector hỗ trợ \cite{cortes1995}, tìm
siêu phẳng có \emph{lề lớn nhất} tới các điểm gần nhất:
\begin{equation}
  \min_{\mathbf{w}, b}\ \tfrac{1}{2}\lVert \mathbf{w} \rVert^2 \quad \text{với}\quad y_i(\mathbf{w}^\top\mathbf{x}_i + b) \ge 1,\ \forall i.
\end{equation}
Nhờ thủ thuật hàm nhân, ví dụ nhân RBF $K(\mathbf{x}, \mathbf{x}') = \exp(-\gamma\lVert\mathbf{x}-\mathbf{x}'\rVert^2)$,
SVM tách được dữ liệu phi tuyến bằng cách ngầm ánh xạ lên không gian vô hạn chiều. Bài toán tối ưu là lồi nên không
mắc cực tiểu địa phương; SVM thống trị nhận dạng ảnh và phân loại văn bản suốt hơn một thập kỷ.

\subsection{Cây quyết định và các phương pháp học tập hợp}

Cây quyết định (CART, C4.5) chia không gian bằng chuỗi câu hỏi đúng/sai, dễ diễn giải. Hai cách kết hợp nhiều mô hình
yếu thành một mô hình mạnh xác lập vị thế trên dữ liệu bảng: \textbf{Boosting} (AdaBoost \cite{freund1997}, sau đó
Gradient Boosting) huấn luyện tuần tự các mô hình sửa lỗi của nhau, còn \textbf{Random Forest} \cite{breiman2001}
lấy trung bình hàng trăm cây học trên các mẫu bootstrap với tập con đặc trưng ngẫu nhiên, giảm phương sai mà không
tăng độ chệch. Chương~\ref{ch:hoc-may} cho thấy Random Forest đến nay vẫn là đối thủ khó vượt trên dữ liệu bảng.

\subsection{LeNet-5 (1998): nền tảng của mạng tích chập hiện đại}

Năm 1998, LeCun và cộng sự công bố LeNet-5 \cite{lecun1998}, ghép ba ý tưởng: vùng tiếp nhận cục bộ, chia sẻ trọng số
và lấy mẫu giảm (pooling). LeNet-5 đọc hàng triệu tấm séc ngân hàng mỗi ngày ở Hoa Kỳ, nhưng giới hạn CPU và thiếu dữ
liệu ảnh lớn khiến nó chưa mở rộng được sang ảnh tự nhiên phức tạp.

\section{Cuộc cách mạng học sâu (2012 -- 2020)}

\subsection{Bước ngoặt ImageNet 2012 và AlexNet}

Tại cuộc thi ILSVRC 2012, AlexNet của Krizhevsky, Sutskever và Hinton \cite{krizhevsky2012} đạt sai số top-5 15,3\,\%,
bỏ xa vị trí thứ hai (26,2\,\%) dùng đặc trưng thủ công SIFT/HOG cùng SVM. Ba yếu tố hội tụ: \textbf{dữ liệu lớn}
(ImageNet hơn 1,2 triệu ảnh, 1.000 lớp), \textbf{GPU} với CUDA tăng tốc nhân ma trận hàng chục lần so với CPU, và
\textbf{các cải tiến thuật toán} như hàm kích hoạt ReLU $f(x) = \max(0, x)$ không bão hòa ở vùng dương và điều chuẩn
Dropout \cite{srivastava2014} ngắt ngẫu nhiên nơ-ron để chống đồng thích nghi.

\subsection{Sự tiến hóa của các kiến trúc thị giác: VGG, ResNet, MobileNet}

\textbf{VGGNet} \cite{simonyan2015} chuẩn hóa thiết kế bằng cách xếp chồng các bộ lọc $3{\times}3$: hai tầng $3{\times}3$
có cùng vùng nhìn với một tầng $5{\times}5$ nhưng ít hơn 28\,\% trọng số và thêm một hàm phi tuyến. \textbf{ResNet}
\cite{he2016} giải quyết hiện tượng thoái hóa (mạng sâu hơn lại có sai số huấn luyện cao hơn) bằng kết nối tắt
$\mathbf{y} = F(\mathbf{x}) + \mathbf{x}$, cho phép huấn luyện mạng 152 tầng và đưa sai số top-5 ImageNet xuống 3,57\,\%.
\textbf{MobileNet} \cite{howard2017,sandler2018} tách tích chập chuẩn thành depthwise và pointwise, giảm chi phí tính
toán 8 -- 9 lần để chạy trên điện thoại. Cả ba kiến trúc được hiện thực và đối chuẩn ở Chương~\ref{ch:cnn}.

\subsection{Mô hình hóa chuỗi: từ RNN, LSTM đến Transformer}

RNN duy trì một trạng thái ẩn qua các bước thời gian, nhưng lan truyền ngược qua thời gian nhân liên tiếp các ma trận
Jacobian nên gradient suy giảm theo hàm mũ \cite{bengio1994}. LSTM \cite{hochreiter1997} thêm ô nhớ được điều khiển bởi
ba cổng (quên, vào, ra), tạo đường đi cộng cho gradient; GRU \cite{cho2014} rút gọn còn hai cổng. Năm 2017, Transformer
\cite{vaswani2017} bỏ hẳn liên kết hồi quy, dùng cơ chế tự chú ý đa đầu để mọi vị trí trong chuỗi nhìn thấy nhau trực
tiếp và tính song song hoàn toàn trên GPU, mở đường cho các mô hình ngôn ngữ khổng lồ. Chương~\ref{ch:rnn} hiện thực
và so sánh bốn kiến trúc hồi quy.

\section{Kỷ nguyên AI tạo sinh, mô hình nền tảng và hệ đa tác tử (2020 -- nay)}

\subsection{Mô hình ngôn ngữ lớn và định luật tỷ lệ}

Kaplan và cộng sự \cite{kaplan2020}, rồi Hoffmann và cộng sự (Chinchilla) \cite{hoffmann2022} cho thấy mất mát của mô
hình ngôn ngữ giảm theo luật lũy thừa khá đều với ba đại lượng: số tham số, lượng dữ liệu và ngân sách tính toán. Ở quy
mô hàng trăm tỷ tham số xuất hiện các năng lực nổi trội: học trong ngữ cảnh từ vài ví dụ mà không cập nhật trọng số,
suy luận từng bước (chain-of-thought), sinh mã nguồn. Hình~\ref{fig:nested} tóm tắt quan hệ bao hàm giữa các thế hệ.

\hinhtikz{f1_nested}{Quan hệ bao hàm giữa trí tuệ nhân tạo, học máy, học sâu và AI tạo sinh / hệ đa tác tử.}{nested}

\subsection{Giới hạn của mô hình thuần tham số và vai trò của đồ thị tri thức, RAG}

Mô hình ngôn ngữ chỉ lưu tri thức trong trọng số nên có ba giới hạn: \emph{ảo giác} (sinh câu trả lời trôi chảy nhưng
sai sự thật, vì mô hình tối ưu xác suất từ tiếp theo chứ không kiểm chứng), \emph{tri thức đóng băng} tại thời điểm
tiền huấn luyện, và \emph{khó truy vết nguồn} trong các lĩnh vực cần trích dẫn như y tế, pháp luật. Xu hướng khắc phục
là ghép thêm bộ nhớ phi tham số: \textbf{đồ thị tri thức} lưu sự kiện dưới dạng bộ ba (thực thể, quan hệ, thực thể) để
suy diễn tường minh; \textbf{mạng nơ-ron đồ thị} mở rộng học sâu sang dữ liệu có cấu trúc đồ thị bằng truyền thông điệp
giữa các nút; \textbf{RAG} \cite{lewis2020} truy xuất tài liệu liên quan rồi đưa vào ngữ cảnh trước khi sinh câu trả lời.

\subsection{Xu hướng: hệ đa tác tử tự trị}

Bước tiếp theo là các \emph{tác tử} (agent) tự chủ có chu trình nhận thức đầy đủ: quan sát, lập kế hoạch và suy luận,
gọi công cụ hay API, rồi tự đánh giá và sửa sai. Nhiều tác tử chuyên biệt (lập kế hoạch, nghiên cứu, lập trình, kiểm thử)
phối hợp và phản biện nhau trong một hệ đa tác tử giải được những việc kỹ thuật mà một mô hình đơn lẻ không làm nổi.
Đây là định hướng then chốt của hệ thống thông minh thế hệ mới, và cũng là nội dung các chương sau của học phần
(đồ thị tri thức, mạng nơ-ron đồ thị, hệ đa tác tử) sẽ kế thừa từ nền tảng mô hình hóa được xây dựng trong tiểu luận này.
